\section{Conclusão}\label{sec:conclusao}

Este trabalho apresentou o Fannar, um método pós-hoc de interpretabilidade mecanicista baseado em \textit{kernels}, que analisa geometricamente fontes internas do modelo. O método associa a cada objeto fontes de informação extraídas do modelo, compara cada fonte por um \textit{kernel} positivamente definido e compõe as matrizes de Gram pelo produto de Hadamard, o que realiza o produto tensorial dos RKHS correspondentes. Dessa composição resultam uma pseudométrica e uma decomposição da energia espectral bem definidas para quaisquer fontes e \textit{kernels}, sobre as quais são formuladas leituras globais, locais e contrastivas. Os resultados teóricos utilizam os fundamentos clássicos de RKHS para garantir que o objeto analisado está bem definido. Se a geometria resultante é informativa sobre o comportamento do modelo é uma questão empírica, respondida para cada instanciação contra referências externas.

Na instanciação para redes neurais, avaliada em dez bases tabulares, no Oxford-IIIT Pet e no CIFAR-100, a geometria reflete a decisão e o erro do modelo acima de um piso construído para cada leitura, a contenção das direções de classe no subespaço retido é a leitura espectral mais robusta, e as distâncias entre classes acompanham a hierarquia semântica das bases de imagens, que o modelo nunca viu. Em todos esses casos, entretanto, a geometria reflete o comportamento sem ir além dele: a margem geométrica não antecipa o erro melhor do que o \textit{softmax}, e a hierarquia das distâncias não excede a que a confusão do modelo já expressa. A composição de fontes acrescenta o gradiente às ativações nos blocos rasos e intermediários das redes convolucionais, e as ablações mostram que o produto de Hadamard e a soma das mesmas fontes produzem as mesmas leituras, enquanto a normalização angular de cada fonte é a escolha decisiva. Na leitura contrastiva, o método supera o SHAP e o LIME na \texttt{letter} e na \texttt{pendigits} e identifica atributos informativos bem acima do acaso, mas perde para as decomposições exatas da margem, o valor de Shapley com o mesmo referencial e os gradientes integrados. Aplicado a uma árvore de decisão, recupera o atributo do nó de separação em 85\% e 91\% dos pares e a profundidade dessa separação em 79\% e 86\%, também sem a fonte que já codifica o nó. As leituras espaciais caracterizam a localização e a organização das representações, enquanto o Grad-CAM caracteriza a sensibilidade da decisão. Os resultados de deleção mostram que são informações distintas, e a leitura espacial não se transfere ao transformador.

A principal contribuição do trabalho é, portanto, menos uma técnica que supera as demais em cada tarefa e mais um arcabouço que separa o que é garantido pela construção do que precisa ser verificado a cada instanciação, acompanhado de um protocolo em que cada leitura tem uma referência de nulidade, uma pergunta sobre o comportamento do modelo e, sempre que possível, uma alternativa simples contra a qual é comparada. Os resultados negativos fazem parte dessa contribuição, visto que delimitam o que a geometria de um modelo pode e não pode revelar sobre ele. Como trabalhos futuros, destacam-se modelos pré-treinados e bases maiores, a leitura de pares por trocas sequenciais, fontes adequadas a transformadores (como a atenção) e menos correlacionadas entre si, e a comparação de modelos distintos em uma mesma geometria.
